\chapter{Recommandations et perspectives}
\label{chap:perspectives}

\section*{Introduction}
\addcontentsline{toc}{section}{Introduction}
Ce chapitre recueille des recommandations méthodologiques pour la
détection multi-étiquettes des émotions et ouvre des perspectives de
prolongement, en lien avec le protocole du
chapitre~\ref{chap:experimentation}.

\section{Principes méthodologiques}
\label{sec:persp-enseignements}

Quatre principes guideront l'évaluation de HAKE-MER et, plus
généralement, tout travail sur GoEmotions :
\begin{itemize}
\item \textbf{Multi-backbone.} L'effet d'un module architectural peut
      dépendre de la capacité de l'encodeur ; conclure à partir d'un
      seul PLM est risqué.
\item \textbf{Contrôle du prétraitement.} Troncature, segmentation en
      phrases et initialisation des têtes peuvent masquer ou imiter
      un gain architectural.
\item \textbf{Métriques par classe.} Le F1-macro seul ne suffit pas :
      les classes rares et la frontière avec \emph{neutral} doivent
      être inspectées explicitement.
\item \textbf{Comparaisons externes prudentes.} Se reporter à la
      littérature n'a de sens qu'avec un protocole comparable ou une
      discussion honnête des écarts.
\end{itemize}

\section{Recommandations pratiques (après expériences)}
\label{sec:persp-recommandations}

Une fois les ablations menées, la configuration retenue devra être
justifiée par les métriques \emph{et} par le coût (temps GPU,
paramètres). Le plan prévoit de comparer M1+M2+M3
avec NRC sur RoBERTa et DistilBERT, puis de tester M4 seulement si
M1--M3 apportent un gain mesurable.

\section{Perspectives de recherche}
\label{sec:persp-recherche}

\begin{itemize}
\item \textbf{Traitement du déséquilibre.} Pertes focales ou
      repondérées, seuils par émotion ;
\item \textbf{Enrichissement sémantique.} Au-delà de NRC/SenticNet,
      graphes de connaissances (ConceptNet, ressources affectives) et
      fusion adaptative par émotion ;
\item \textbf{Validation externe.} SemEval-2018 Tâche~1, plusieurs
      graines aléatoires ;
\item \textbf{Émotions implicites.} Ressources de bon sens (ATOMIC,
      COMET) ;
\item \textbf{LLM.} Positionner HAKE-MER par rapport aux modèles
      généralistes et aux LLM affectifs spécialisés lorsque des
      baselines reproductibles sont disponibles.
\end{itemize}

\section*{Conclusion}
\addcontentsline{toc}{section}{Conclusion}
Les perspectives ci-dessus prolongent le sujet au-delà du protocole
détaillé au chapitre~\ref{chap:experimentation}.
